%%%%%%%%%%%%%%%%%%%%%%%%%%%%%%%%%%%%%%%%%
% MUW Poster
% LaTeX Template
% Version 1.0 (31/08/2016)
% (Based on Version 1.0 (31/08/2015) of the Jacobs Portrait Poster
%
% License:
% CC BY-NC-SA 3.0 (http://creativecommons.org/licenses/by-nc-sa/3.0/)
%
% Created by:
% Nicolas Ballarini, CeMSIIS, Medical University of Vienna
% nicoballarini@gmail.com
% http://statistics.msi.meduniwien.ac.at/
%
%
% Adaptado por Augusto Ferreira, Universidade Federal de Lavras
% gstfrr@gmail.com
%
%%%%%%%%%%%%%%%%%%%%%%%%%%%%%%%%%%%%%%%%%


\def\footer#1{\def\insertfooter{#1}}
%-----------------------------------------------
%	PACOTES E CONFIGURAÇÕES DO DOCUMENTO
%-----------------------------------------------

\documentclass[final]{beamer}
\usepackage[brazil]{babel}
%\usepackage[utf8]{inputenc}
\usepackage[T1]{fontenc}

\usepackage[scale=1.1]{beamerposter}
\usetheme{MUWposter} % Use the MUWposter theme supplied with this template

\logo{
	\pgfputat{\pgfxy(-16,82)}{
		\pgfbox[center,base]{
			% Se seu trabalho possuir uma logo, coloque-a aqui
			\includegraphics[width=10cm]{ProjectLogo.png}
		}
	}
}

\usepackage{multicol}
\usepackage{array}
\usepackage{pgf}  
\usepackage{float}
\usepackage{mathtools}
\usepackage{amsmath, amsthm, amssymb, amsfonts}
\usepackage{exscale}
\usepackage{xcolor}
\usepackage{ushort}
\usepackage{setspace}
\usepackage{url}
\usefonttheme{default} 

\definecolor{azulERMAC}{RGB}{82,169,209}

%Pasta das imagens
\graphicspath{{./imagens/}}

% Pacotes de referências bibliográficas
\usepackage[alf]{abntex2cite}

%-----------------------------------------------
%  CORES E PLANO DE FUNDO
%-----------------------------------------------
\colorlet{themecolor}{azulERMAC}
\usebackgroundtemplate{\includegraphics[width=\paperwidth]{fundo_ermac.jpg}}
%-----------------------------------------------
%  FIM CORES E PLANO DE FUNDO
%-----------------------------------------------


%-----------------------------------------------
%  INÍCIO Largura das colunas
%-----------------------------------------------
\setlength{\paperwidth}{85cm} % Largura do poster
\setlength{\paperheight}{100cm} % Altura do poster
\newlength{\sepmargin}
\newlength{\sepwid}
\newlength{\onecolwid}
\newlength{\twocolwid}
\newlength{\threecolwid}

% Configurações das colunas
\setlength{\sepmargin}{0.055\paperwidth} % Separation width (white space) between columns
\setlength{\sepwid}{0.03\paperwidth} % Separation width (white space) between columns
\setlength{\onecolwid}{0.43\paperwidth} % Largura de uma coluna
\setlength{\twocolwid}{0.9\paperwidth} % Largura de duas colunas

%-----------------------------------------------
%  FIM Largura das colunas
%-----------------------------------------------


%--------------------------------------------------------------------------------------
%	SEÇÃO DO TÍTULO
%--------------------------------------------------------------------------------------
\setbeamertemplate{title}[left]
\setbeamertemplate{frametitle}[default][left]

\title{Título do trabalho} % Título do poster

\author{Primeiro Autor, Segundo Autor} % Autor(es)

\institute{Instituição do Primeiro Autor, Instituição do Segundo Autor} % Instituições
% Os emails ficarão na parte de Contato, lá no final
%--------------------------------------------------------------------------------------



\begin{document}

\addtobeamertemplate{block end}{}{\vspace*{1ex}} % Espaço após os blocos
\addtobeamertemplate{block alerted end}{}{\vspace*{0ex}} % Espaço do bloco de contato
\setlength{\belowcaptionskip}{2ex} % Espaço embaixo de figuras
\setlength\belowdisplayshortskip{1ex} %Espaço embaixo de equações

\begin{frame}[t] % O poster todo é um frame do beamer (gambiarra?? talvez)

	\begin{columns}[t] % O conteúdo do poster será dividido entre duas colunas

		\begin{column}{\sepmargin}
		\end{column}


		%-----------------------------------------------
		%  Primeira Coluna
		%-----------------------------------------------

		\begin{column}{\onecolwid} % Primeira Coluna

			\begin{block}{Introdução} % Cada seção de texto das colunas é um bloco
				Uptam, offictibus rem vendipici nonsecest la cullabor moluptas exero quo blam quamentum repudistiam iunda doluptate dolorecta quatem faceaquodit optatum nonse init volori doluptas nam erferch ilique comnihil ma doluptate sanditat ommo temquia nonse sed modicium que vollacillab ius. Uptam, offictibus rem vendipici nonsecest la cullabor moluptas exero quo blam.
			\end{block}

			\begin{block}{Objetivos}
				Uptam, offictibus rem vendipici nonsecest la cullabor moluptas exero quo blam quamentum repudistiam iunda doluptate dolorecta quatem faceaquodit optatum nonse init volori doluptas nam erferch ilique.
			\end{block}

			\begin{block}{Metodologia}
				Uptam, offictibus rem vendipici nonsecest la cullabor moluptas exero quo blam quamentum repudistiam iunda doluptate dolorecta quatem faceaquodit optatum nonse init volori doluptas nam erferch ilique comnihil ma doluptate sanditat ommo temquia nonse sed modicium que vollacillab ius. Uptam, offictibus rem vendipici nonsecest la cullabor moluptas exero quo blam.
			\end{block}

			\begin{table}[ht] % Modelo de tabela
				\centering

				\begin{tabular}{lccc}\hline
					Patient Characteristics & n 		& & \%  \\
					\hline
					Total of Patients 		& 100     	& & 100 \\
					Age 	 				&  45      	& &  45 \\
					Gender 					&  50 	   	& &  50 \\
					\hline
				\end{tabular}
				\caption{Considered scenarios} 
				\label{Patient and tumor Characteristics}
			\end{table}


			\begin{block}{Resultado}
				Uptam, offictibus rem vendipici nonsecest la cullabor moluptas exero quo blam quamentum repudistiam iunda doluptate dolorecta quatem faceaquodit optatum nonse init volori doluptas nam erferch ilique comnihil ma doluptate sanditat ommo temquia nonse sed modicium que vollacillab ius. Uptam, offictibus rem vendipici nonsecest la cullabor moluptas exero quo blam.
				\begin{equation}
					\int\frac{\sqrt[a]{xˆb}}{bˆ{ax}}dx=\gamma F(a)+ \delta \pi F(bn)
				\end{equation}

				Uptam, repudistiam iunda doluptate dolorecta quatem faceaquodit optatum nonse init volori doluptas nam erferch ilique comnihil ma doluptate sanditat ommo temquia nonse sed modicium que vollacillab ius. Uptam, offictibus rem vendipici nonsecest la cullabor moluptas exero quo blam.
			\end{block}

			\begin{block}{Conclusão}
				Uptam, offictibus rem vendipici nonsecest la cullabor moluptas exero quo blam quamentum repudistiam iunda doluptate dolorecta quatem faceaquodit optatum.
			\end{block}

		\end{column}


		\begin{column}{\sepwid}
		\end{column}


		%-----------------------------------------------
		%  Segunda Coluna
		%-----------------------------------------------
		\begin{column}{\onecolwid} % Segunda Coluna

			\begin{block}{\vspace*{2.7cm}}
				Uptam, offictibus rem vendipici nonsecest la cullabor moluptas exero quo blam quamentum repudistiam iunda doluptate dolorecta quatem faceaquodit optatum nonse init volori doluptas nam erferch ilique comnihil ma doluptate sanditat ommo temquia nonse sed modicium que vollacillab ius. Uptam, offictibus rem vendipici nonsecest la cullabor moluptas exero quo blam \citeonline{agudo1964introduccao}.
				Nonse init volori doluptas nam erferch ilique comnihil ma doluptate sanditat ommo temquia nonse sed modicium que vollacillab ius. Uptam, offictibus rem vendipici nonsecest la cullabor moluptas exero quo blam.
			\end{block}

			\begin{block}{}

				\begin{figure} % Exemplo de imagem
					\vspace*{-1cm}
					\includegraphics[width=.7\linewidth]{plot.PNG}
				\end{figure}
				
				\begin{figure}
					\includegraphics[width=.7\linewidth]{Bild1.jpg}
				\end{figure}

				\begin{multicols}{2} % Também é possível colocar duas imagens lado a lado
					\begin{figure}
						\vspace*{-0.95cm}
						\includegraphics[width=.5\linewidth]{Bild2.jpg}
					\end{figure}
					
					\begin{figure}
						\vspace*{-0.95cm}
						\includegraphics[width=.5\linewidth]{Bild3.jpg}
					\end{figure}
				\end{multicols}

			\end{block}
		\end{column}

		\begin{column}{\sepmargin} \end{column}
	\end{columns} % Fim das colunas de conteúdo


	\begin{center}
	\vspace{5pt}
	\noindent\textcolor{azulERMAC}{\rule{75cm}{0.8pt}}
	\vspace{-5pt}
	\end{center}

	
	%-----------------------------------------------
	%	PARTE DE BAIXO
	%-----------------------------------------------

	\begin{columns}[b] % Criação de duas colunas

		\begin{column}{\sepmargin}
		\end{column}

		\begin{column}{\onecolwid} % Contato

			\vspace{0pt}

			\begin{alertblock}{\large Contato}
				%\vspace*{-0.5cm}
				\begin{footnotesize}
					\begin{itemize}
						\item \href{mailto:email@meduniwien.ac.at}{email@meduniwien.ac.at}
						\item \href{http://www.example.com/}{www.example.com} - \href{www.meduniwien.ac.at/medstat}{www.meduniwien.ac.at/medstat}
					\end{itemize}
				\end{footnotesize}	
			\end{alertblock}

		
		\end{column}% Fim da coluna Contato

		\begin{column}{\sepwid}
		\end{column} %

		\begin{column}{\onecolwid} % Referências

			%-----------------------------------------             
			%       REFERÊNCIAS BIBLIOGRÁFICAS      
			%-----------------------------------------        
			\begin{block}{\large Referências}
				\vspace{0pt}
				\nocite{meneghessocodigos,arantes2007,rmc_wiki} % opcional para referenciar quem não foi citado
				{\footnotesize
				\bibliography{bibliog.bib}}
			\end{block} 
		\vspace{0pt}
		\end{column} % Fim da coluna Referências

		\begin{column}{\sepmargin}
		\end{column}

	\end{columns} % Final das colunas da parte de baixo

\end{frame} %Fim do frame principal
\end{document}